\documentclass[11pt]{article}
\usepackage[T1]{fontenc}
\usepackage[utf8]{inputenc}
\usepackage{lmodern}
\usepackage[margin=1in]{geometry}
\usepackage{amsmath,amssymb,longtable,array,calc}
\usepackage[colorlinks=true,linkcolor=blue,urlcolor=blue]{hyperref}
\IfFileExists{booktabs.sty}{\usepackage{booktabs}}{%
  \newcommand{\toprule}{\hline}
  \newcommand{\midrule}{\hline}
  \newcommand{\bottomrule}{\hline}
}
\setlength{\emergencystretch}{3em}
\setlength{\parindent}{0pt}
\setlength{\parskip}{5pt}
\allowdisplaybreaks
\providecommand{\tightlist}{\setlength{\itemsep}{0pt}\setlength{\parskip}{0pt}}
\newcounter{none}
\title{An Edge-Conditioning Reduction for Perfect-Matching Eigenvalues}
\author{Claude Make No Mistakes}
\date{25 November 2026 --- Working paper}

\begin{document}
\maketitle
\begin{abstract}
Conditioning on a shared edge reduces the signed second-eigenvalue
inequality for each fixed core to the case of exactly two shared edges.
We prove the reduction and verify the bases for 230 cores of weight at
most sixteen. A separate argument settles the single-cycle family for
arbitrary cycle size.
We also prove the minimum-gap conjecture for all nonidentity relations at
every size $n\ge5$ and show that the two-edge switch uniquely attains the
minimum.
\end{abstract}

For a fixed nonempty partition \(\nu\vdash s\) with parts at least two,
the signed second-eigenvalue conjecture for \((\nu,1^{n-s})\) at all
\(n\ge s+2\) is equivalent to its single instance at \(n=s+2\). When the
target attains the signed maximum at the base, at each later size it is
strictly above every competing eigenvalue on a different nontrivial
representation. Exact orbit quotients settle the four bases of size
\(s=6\) at \(n=8\), which proves the conjecture for all four families at
every \(n\ge8\).

\section{Definitions and the signed
target}\label{definitions-and-the-signed-target}

Write \(\Omega_n\) for the perfect matchings of a complete graph on
\(2n\) vertices. The relation type of \((M,N)\) records the half-lengths
of the alternating cycles of the two matchings, counting a shared edge
as a doubled cycle of half-length one. Fix \(\nu\) as above and put
\(\mu_n=(\nu,1^{n-s})\) for \(n\ge s\). Let \(A_n\) be its adjacency
operator, \(d_n\) its valency, and \(P_n=A_n/d_n\). Write \(b_n\) for
the largest eigenvalue of \(P_n\) on the orthogonal complement of
constants. This is the second eigenvalue in decreasing signed order with
multiplicities. All inner products on matching spaces below use uniform
measure.

The matching permutation module is the multiplicity-free sum of the
Specht modules \(V_{2\lambda}\), for \(\lambda\vdash n\). In the
authors' notation the target row is \(\lambda=(n-1,1)\), and the
subscript of \(A\) is the relation, not the module. Their formal
Conjecture 1.3 also covers the one-singleton hook \((n-1,1)\) except at
\(n=4\). Their Corollary 3.5 handles those hooks for \(n\ge6\), and
Appendix A handles \(n=3,5\).
\href{https://arxiv.org/html/2510.17135v2}{Gupta et al., Conjecture 1.3,
Section 2 and Corollary 3.5}.

For an edge \(e\) of the complete graph let \(g_e(M)=\mathbf1_{e\in M}\)
and let \(U_n=\operatorname{span}\{g_e\}\). The incidence map from the
permutation module on two-element sets to the matching module is
equivariant. By Young's rule the domain has constituents
\(V_{(2n)},V_{(2n-1,1)},V_{(2n-2,2)}\) for \(n\ge2\). The target
permutation module contains no \(V_{(2n-1,1)}\) because its constituents
have even row lengths. Constants lie in the image since
\(\sum_e g_e=n\), and \(g_e\) is nonconstant for \(n\ge2\). Consequently

\[
U_n=\mathbf R\mathbf1\oplus V_{(2n-2,2)}.
\]

Conditional on \(M\), a uniform \(\mu_n\)-neighbor \(N\) shares exactly
\(n-s\) of the \(n\) edges of \(M\). Its marginal probability of
containing \(e\in M\) is \(a=(n-s)/n\) by symmetry. The stabilizer of
\(M\) is also transitive on the \(\binom{2n}{2}-n=2n(n-1)\) edges
outside \(M\), so the marginal there is \(c=s/[2n(n-1)]\). Thus
\(P_ng_e=c+(a-c)g_e\), and subtracting the mean \(1/(2n-1)\) gives the
normalized target eigenvalue

\[
\rho_n=a-c=1-\frac{s}{2}\left(\frac1n+\frac1{n-1}\right).
\]

The valency is the number of ways to select the \(s\) matching edges to
change times the valency on those \(2s\) vertices. In symbols,
\(d_n=\binom ns d_s\) and \(d_s=2^{s-\ell(\nu)}s!/\prod_i i^{m_i}m_i!\),
where \(m_i\) is the multiplicity of \(i\) in \(\nu\). The target
spectral gap is \(d_n(1-\rho_n)\) wherever the conjecture holds.
\href{https://arxiv.org/html/2510.17135v2\#S2.SS1}{Gupta et al., Lemmas
2.1--2.2}.

\section{Restriction to a shared
edge}\label{restriction-to-a-shared-edge}

Suppose \(n>s\). For any edge \(e\), let \(\Omega_e=\{M:e\in M\}\) and
let \(f_e=f|_{\Omega_e}\). Removing the endpoints of \(e\) identifies
\(\Omega_e\) with \(\Omega_{n-1}\). The induced relation is
\(\mu_{n-1}\), and its normalized adjacency is \(P_{n-1}\) after this
identification.

Consider the following sampling rule. Choose \(e\) uniformly among all
edges of the complete graph, then choose \(M\) uniformly in
\(\Omega_e\), then choose \(N\) uniformly among its induced neighbors in
\(\Omega_e\). Every matching contains exactly \(n\) edges and all stars
have the same size, so \(M\) is uniform and \(e\) given \(M\) is uniform
among its \(n\) edges. For a fixed \(\mu_n\)-neighbor \(N\), there are
exactly \(n-s\) eligible choices of \(e\in M\cap N\). The conditional
probability of that \(N\) given \(M\) is \((n-s)/(nd_{n-1})=1/d_n\).
This proves the exact identities

\[
\langle f,P_nf\rangle=\mathbb E_e\langle f_e,P_{n-1}f_e\rangle_{\Omega_e},
\qquad
\|f\|^2=\mathbb E_e\|f_e\|^2_{\Omega_e}.
\]

Take \(f\in W_n=U_n^\perp\). For every edge \(e\), orthogonality to
\(g_e\) says that \(f_e\) has mean zero. The local signed Rayleigh bound
yields

\[
\langle f,P_nf\rangle\le b_{n-1}\|f\|^2.
\]

Both \(U_n\) and \(W_n\) are invariant. The target row gives
\(b_n\ge\rho_n\) and the Rayleigh inequality gives
\(b_n\le\max(\rho_n,b_{n-1})\). For every \(n>s\) in range,
\(\rho_n-\rho_{n-1}=s/[n(n-2)]>0\). If \(b_{n-1}\le\rho_n\), the
conjectured row attains \(b_n\). If \(b_{n-1}<\rho_n\), the top
nonconstant eigenspace is exactly \(V_{(2n-2,2)}\). In particular,
starting from \(b_{s+2}=\rho_{s+2}\), even with a tie at the base,
proves the result and removes any tie across modules at every later
size. The converse is immediate by taking \(n=s+2\). Since \(\rho_n<1\)
for a nonempty core, any size certified this way also has a connected
relation graph.

This argument requires at least one shared edge for a restriction step.
The identity relation has \(s=0\) and every eigenvalue is 1. For \(n=1\)
the space has dimension one and the second-eigenvalue target has no
meaning. The one-singleton relation \((3,1)\) at \(n=4\) has target 4
and actual second eigenvalue 8 on \(V_{(2,2,2,2)}\). Even outside hooks,
the example \((2,2,1)\) at \(n=5\) has target 6 and second eigenvalue
15. A signed inequality can hold when the absolute-value version fails.
At \((n,\mu)=(4,(2,1,1))\) the target is 5, and one competing eigenvalue
is \(-6\). \href{https://arxiv.org/html/2510.17135v2\#A1}{Gupta et al.,
Appendix A, Tables 5--6}.

\subsection{A corollary with unbounded core
size}\label{a-corollary-with-unbounded-core-size}

For every pair of integers \(s\ge2\) and \(n\ge s+2\), the relation
\((s,1^{n-s})\) has its largest nonconstant eigenvalue on
\(V_{(2n-2,2)}\). For \(s\ge5\), start from the published Corollary 3.5
at size \(s+1\ge6\). It gives \(b_{s+1}=\rho_{s+1}\) for the
one-singleton relation \((s,1)\). Since \(\rho_{s+1}<\rho_{s+2}\), the
preceding argument applies at every size from \(s+2\) onward. For
\(s=2,3,4\), use the exact two-singleton bases at \(n=4,5,6\) in the
computed table. Both the size of the changed alternating cycle and the
number of shared edges can therefore vary without bound.
\href{https://arxiv.org/html/2510.17135v2\#S3}{Gupta et al., Corollary
3.5}.

\section{Exact bases for the weight-six
families}\label{exact-bases-for-the-weight-six-families}

Use the adjacent-pair matching \(M_0\) and let \(H_n\) be its
stabilizer. Its orbits on \(\Omega_n\) are indexed by the relation type
with \(M_0\). To see transitivity for a given type, match up alternating
cycles of the same length, then relabel corresponding vertices around
each cycle, preserving the \(M_0\)-edges. For each relation \(\mu\),
orbit row \(\sigma\), and orbit column \(\tau\), the quotient entry is
the number of \(N\) with pair types
\(\operatorname{type}(M_\sigma,N)=\mu\) and
\(\operatorname{type}(M_0,N)=\tau\), where \(M_\sigma\) is any
representative. These quotients act on the \(H_n\)-fixed subspace. Every
matching-module constituent has a one-dimensional fixed subspace by
Frobenius reciprocity and the multiplicity-free decomposition. Thus the
quotient sees every constituent once, including constituents that tie
for a particular relation.
\href{https://doi.org/10.5802/alco.104}{Srinivasan, The perfect matching
association scheme}.

The native enumerator accumulated all the counts from 2,027,025
matchings at \(n=8\). It enumerates each matching by pairing the least
remaining vertex with each available partner in turn. A walk around the
alternating cycles gives the relation types. It stores the 22 quotients
of size \(22\times22\) in integer counters. The Python implementation
computes the same quotients separately for \(n=3,4\) and verifies the
complete published spectra at those sizes. The checks at all three sizes
include the diagonal relation, row sums, detailed balance with orbit
sizes, commutativity, and exact joint eigenvectors over the rationals.

The joint eigenvectors were labeled independently of the workbook. If
\(u\) is the spherical vector scaled so that \(u(M_0)=1\), its uniform
squared norm is \(1/\dim V_{2\lambda}\). Indeed, if \(Q\) is the
orthogonal projector onto the constituent, then \(Q\delta_{M_0}\) is
fixed by \(H_n\), its squared counting norm is \(Q_{M_0,M_0}\), and
transitivity gives \(Q_{M_0,M_0}=\dim V_{2\lambda}/|\Omega_n|\). The
central transposition sum contributes \(n\) copies of the original
matching and reaches each switch-neighbor twice, so its action is
\(nI+2A_{(2,1^{n-2})}\). Its scalar on \(V_{2\lambda}\) is the sum of
the Young diagram contents. The pair of hook-length dimension and
content sum separates the labels at these sizes.
\href{https://arxiv.org/html/2510.17135v2\#S4}{Gupta et al., Table 1 and
Proposition 4.1}.

{\def\LTcaptype{none} % do not increment counter
\begin{longtable}[]{@{}
  >{\raggedright\arraybackslash}p{(\linewidth - 6\tabcolsep) * \real{0.2000}}
  >{\raggedleft\arraybackslash}p{(\linewidth - 6\tabcolsep) * \real{0.2667}}
  >{\raggedleft\arraybackslash}p{(\linewidth - 6\tabcolsep) * \real{0.2667}}
  >{\raggedleft\arraybackslash}p{(\linewidth - 6\tabcolsep) * \real{0.2667}}@{}}
\toprule\noalign{}
\begin{minipage}[b]{\linewidth}\raggedright
Base relation
\end{minipage} & \begin{minipage}[b]{\linewidth}\raggedleft
Valency
\end{minipage} & \begin{minipage}[b]{\linewidth}\raggedleft
Target on \(V_{(14,2)}\)
\end{minipage} & \begin{minipage}[b]{\linewidth}\raggedleft
Largest other nonconstant eigenvalue
\end{minipage} \\
\midrule\noalign{}
\endhead
\bottomrule\noalign{}
\endlastfoot
\((6,1,1)\) & 107520 & 21120 & 960 \\
\((4,2,1,1)\) & 40320 & 7920 & 2520 \\
\((3,3,1,1)\) & 17920 & 3520 & 1504 \\
\((2,2,2,1,1)\) & 3360 & 660 & 396 \\
\end{longtable}
}

The normalized target is \(11/56\) in each case. The full exact matrix
and both shape orders are in certificate\_n8.json. All 484 adjacency
eigenvalues and all 22 dimensions agree with the pinned authors'
workbook. The worksheet rows use doubled Specht labels and the columns
use doubled relation partitions. Worksheet row 27 finds the first signed
maximum after excluding the trivial row, so ties select the first
matching row. The workbook's per-column test row applies to relations
with at least two doubled parts equal to 2. The Sage function
\texttt{main\_conjecture} uses the same signed maximum and scope and
excludes the identity column. These automated checks do not include the
formal conjecture's one-singleton hooks. The exact comparison above
reads the integer values in the spectrum table and independently
determines all maximizers.
\href{https://github.com/Himanshugupta23/Perfect-Matching-Association-Scheme/blob/9bfd0192a04b1c271b305ee3c8146b356ac41603/Excel_format_Evals_of_PM_Schm_upto_n\%3D15.xlsx}{Authors'
workbook},
\href{https://github.com/Himanshugupta23/Perfect-Matching-Association-Scheme/blob/9bfd0192a04b1c271b305ee3c8146b356ac41603/Evals_of_PM_scheme_via_Srinivasan_algorithm.sage.py\#L245-L274}{authors'
Sage code}.

A separate exact calculation using Jack polynomials certifies all 230
bases for \(2\le s\le16\) and checks 72,035 workbook entries. The same
propagation argument therefore settles every one of those core families
at all admissible sizes. \href{https://github.com/claude-make-no-mistakes/perfect-matching-eigenvalues/blob/main/verification/ZONAL-SPEC.md}{Jack calculation}, \href{https://github.com/claude-make-no-mistakes/perfect-matching-eigenvalues/blob/main/verification/two-singleton-bases.json}{base certificate}.

\section{A lower-order sufficient
condition}\label{a-lower-order-sufficient-condition}

The same Rayleigh identity applies at \(n=s+2\) before any conjecture is
known for one-singleton relations. It suffices to show, for each core,
that

\[
b_{s+1}(\nu,1)\le\rho_{s+2}=\frac{3s+4}{2(s+1)(s+2)}.
\]

This requires only the largest signed nonconstant eigenvalue, regardless
of its label. Exact rational checks of the pinned workbook satisfy the
strict inequality for all 134 cores with \(2\le s\le14\). Iterating once
gives another sufficient condition, \(b_s(\nu)\le\rho_{s+2}\), because
\(\rho_{s+1}<\rho_{s+2}\). That stronger condition fails for \((2,2)\),
\((2,2,2)\), \((3,3)\), and \((2,2,2,2)\) in this range. For example
\(b_4(2,2)=7/12>4/15=\rho_6\). The failures concern the sufficient
bound. The corresponding two-singleton bases still hold. The universal
one-singleton bound and arbitrary two-singleton bases remain open here.

For a shape-specific version, let \(f\) lie in \(V_{2\lambda}\). Under
restriction to the subgroup fixing the endpoints of \(e\) pointwise, the
only even constituents have shapes \(2\kappa\) obtained by removing a
corner of \(\lambda\) and doubling row lengths. This follows by applying
the symmetric-group branching rule twice, since reaching another even
shape requires both removals from the same row. The restriction map
\(f\mapsto f_e\) is equivariant for that subgroup, so the projections of
\(f_e\) lie only in these constituents of the smaller matching module.
The projected squared norms, averaged over \(e\) and divided by
\(\|f\|^2\), are nonnegative and sum to 1. The exact Rayleigh identity
writes the eigenvalue at \(\lambda\) as their convex combination of the
corresponding predecessor eigenvalues. This restricts which lower-order
rows can contribute to a particular competitor. The argument still
requires a bound on the relevant lower-order spectra.

\section{The smallest spectral gap in the perfect matching association
scheme}\label{the-smallest-spectral-gap-in-the-perfect-matching-association-scheme}

\textbf{Theorem.} Let \(n\ge5\). Among nonidentity relations of the
perfect matching association scheme on \(2n\) vertices, the minimum
unnormalized adjacency spectral gap is \(2n-1\). Only the relation
\(\sigma=(2,1^{n-2})\), which switches two edges of a matching, attains
it. This proves Conjecture 6.1 of the arXiv version of Gupta, Herman,
Lacaze-Masmonteil, Maleki and Meagher.
\href{https://arxiv.org/html/2510.17135v2\#S6}{Source, §6}

Write \(\Omega\) for the perfect matchings. The relation type of two
matchings is the partition formed by the half-lengths of their
alternating cycles, with a shared edge contributing a part of size one.
Write \(A_\mu,d_\mu\) for the adjacency matrix and degree of the
relation graph, and define its gap \(\Delta_\mu\) by the smallest
Rayleigh quotient of \(d_\mu I-A_\mu\) on the orthogonal complement of
constants. This definition also gives zero if the graph is disconnected.

\subsection{Comparing each relation to a
switch}\label{comparing-each-relation-to-a-switch}

Take a pair \(M,L\) of the given nonidentity type \(\mu\). Select two
edges of \(M\) in a nontrivial alternating cycle of \(M\cup L\), and
delete them. The cycle becomes two paths whose endpoints are the four
endpoints of the deleted edges. There are exactly three pairings of
those endpoints. Connecting the endpoints within each path makes two
alternating cycles. Each of the other two pairings joins the paths into
one alternating cycle of the original length. One is the original
pairing. Use the other to obtain a matching \(N\). The types of
\((M,L),(L,N),(M,N)\) are then \(\mu,\mu,\sigma\). This construction
also works on a cycle of length four, where the two paths each consist
of a single edge of \(L\).

Relabel this triple using a uniform permutation of the \(2n\) vertices.
For any fixed type, the symmetric group acts transitively on ordered
pairs of that type, by matching their alternating cycles and preserving
the two edge colors. Hence the relabeling gives the uniform ordered pair
law of type \(\mu\) on each leg and the uniform ordered pair law of type
\(\sigma\) on the endpoints.

For \(f\colon\Omega\to\mathbb R\), let \(D_\tau(f)\) denote the mean of
\((f(P)-f(Q))^2\) over uniform ordered pairs of type \(\tau\). The
square inequality \((u+v)^2\le 2(u^2+v^2)\) on the relabeled triple
yields

\[
D_\sigma(f)\le4D_\mu(f).
\]

Under the uniform inner product,
\(\langle f,(d_\tau I-A_\tau)f\rangle=d_\tau D_\tau(f)/2\). The known
switch degree and gap are \(d_\sigma=n(n-1)\) and
\(\Delta_\sigma=2n-1\). Applying the Rayleigh inequality to every \(f\)
perpendicular to constants now gives

\[
\Delta_\mu\ \ge\ \frac{d_\mu}{4n(n-1)}(2n-1).
\]

In particular, the gap is positive for every nonidentity type at these
sizes. \href{https://arxiv.org/html/2510.17135v2\#S4}{Switch eigenvalue
and gap, Gupta et al., §4}

The comparison is strict for any nonconstant \(f\). Equality in the
averaged square inequality would force \(f(L)=(f(M)+f(N))/2\) on the
entire orbit. The three vertex marginals are uniform. Therefore the mean
square of \(f(L)\) equals that of \(f(M)\), while the mean square of
\((f(M)+f(N))/2\) is the mean square of \(f(M)\) minus
\(D_\sigma(f)/4\). This would force \(D_\sigma(f)=0\). The switch gap is
positive, so \(D_\sigma(f)=0\) implies \(f\) is constant. Applying
strictness to a nonzero minimizing eigenvector perpendicular to
constants makes the gap bound strict.

\subsection{Degrees and the small
exception}\label{degrees-and-the-small-exception}

Delete the singleton parts of \(\mu\) to get a core \(\nu\) of weight
\(s\). Then \(d_\mu=\binom{n}{s}d_\nu\), since there are
\(\binom{n}{s}\) ways to choose the reference edges in the nontrivial
components. For a component of half-length \(a\) on prescribed reference
edges the degree is \(2^{a-1}(a-1)!\), by the valency formula for a
single cycle. Fix one of the \(s\) reference edges. Choosing an
\(a\)-component containing it gives at least
\(\binom{s-1}{a-1}2^{a-1}(a-1)!\) arrangements, multiplied by the degree
of the remaining core. The chosen component is uniquely identifiable
from the resulting matching as the component containing the fixed edge.
Count an empty remaining core with degree one.
\href{https://arxiv.org/html/2510.17135v2\#S2.SS1}{Valency formula,
Gupta et al., Lemma 2.1}

Suppose \(s\ge5\), and take \(a\) to be the largest part of the core.

\begin{itemize}
\tightlist
\item
  If \(a\ge4\), the count is at least \(8(s-1)(s-2)(s-3)\). Since
  \(2(s-2)(s-3)\ge s\) for \(s\ge5\), this is at least \(4s(s-1)\).
\item
  If \(a=3\), the residual core has at least one nontrivial component,
  so its degree is at least two. The count is at least \(8(s-1)(s-2)\),
  which is at least \(4s(s-1)\).
\item
  If \(a=2\), all parts are two. For even \(s\ge6\), the degrees \(d_s\)
  satisfy \(d_s=2(s-1)d_{s-2}\), by choosing the partner of the fixed
  reference edge and one of the two possible switches. The base is
  \(d_6=2\cdot5\cdot12=120=4\cdot6\cdot5\). The recurrence preserves the
  bound \(d_s\ge4s(s-1)\), since \(2(s-2)(s-3)\ge s\).
\end{itemize}

Thus \(d_\mu\ge4n(n-1)\) for \(s\ge5\), as
\(\binom ns s(s-1)=n(n-1)\binom{n-2}{s-2}\). For \(s=3\) the degree is
\(8\binom{n}{3}\), which meets the same bound at \(n\ge5\). For \(s=4\)
the two possible core degrees are \(48\) and \(12\), for \((4)\) and
\((2,2)\) respectively. Both meet the bound at \(n\ge6\), and at \(n=5\)
only \(\mu=(2,2,1)\) falls below it. The published \(n=5\) eigenmatrix
gives that relation valency 60 and largest nonconstant eigenvalue 15,
for a gap of \(45>9\). The core of weight two is precisely the switch
relation. These cases establish the minimum.
\href{https://arxiv.org/html/2510.17135v2\#A1}{Exact \(n=5\) table}

\subsection{Scope}\label{scope}

The argument applies to the unnormalized adjacency gap in Conjecture
6.1. The normalized gaps satisfy the factor-four comparison above. The
proof uses a two-step walk for every switch, which by itself yields only
the diameter bound \(2(n-1)\). Conjecture 6.2 asks for \(n-1\). The
separate finite calculation checks that sharper bound over its stated
range.


The lower limit \(n\ge5\) is sharp. At \(n=4\), the \((2,2)\)
relation has gap \(12-7=5\), while the switch relation has gap
\(12-5=7\), by the published \(n=4\) eigenmatrix.
\href{https://arxiv.org/html/2510.17135v2\#A1}{Source, Table 5}

\section{Computational
reproducibility}\label{computational-reproducibility}

The companion repository includes replay commands, exact input data, and
source code. AI assistance was used to develop and check the arguments
and the software. This working paper reports the stated restricted
results.

\end{document}
